\documentclass{article}
\usepackage[utf8]{inputenc}
\usepackage[T1]{fontenc}
\usepackage[margin=2cm]{geometry}
\usepackage{hyperref}
\sloppy
\begin{document}
\section*{Historical source}
Encyclopaedia Britannica scan, source row 604. Parse 3.5 extracted the text. Super 3.5 VL corrected it from an enlarged crop of each block. It can still contain recognition errors.
Page headers, footers, and page numbers are removed.

\subsection*{Untitled continuation (E0055)}
\textit{Continues an entry from an earlier page.}

Thos. attended, walking as mourners. The queen and princesses also assisted at the ceremony, and the concourse of people, including the public bodies, municipal and academical, amounted to many thousands. An interesting part of the ceremony was the performance of a requiem written by the sculptor's intimate and dear friend Ehlenschlager.

Thorwaldsen was rather above the middle height, firmly built, with broad shoulders. The outline of his face was rather square. His general expression was calm and thoughtful, and very pleasing. His eyes were light and penetrating; his mouth wide and usually closely shut. He wore his hair, which had grown nearly white, in large masses divided over his fine broad forehead, and falling on each side, something like a lion's mane, nearly to his shoulders.

something like a non's mane, nearly to his shoulders.
Thorwaldsen will always take rank with Canova and
Flaxman as one of the great regenerators of Sculpture.
The former, his elders, were almost the first to break
through the miserable routine of the worst mixed Italian
and French art which had descended from Bernini and his
school. Thorwaldsen nobly aided in advancing this revolu-
tion, and with them assisted in establishing the practice of
sculpture upon true principles. (r. w.—r.)

\subsection*{THOU, JACQUES AUGUSTE DE (E0056)}
THOU, JACQUES AUGUSTE DE, or, as his name is frequently written in its Latinized form, JACOBUS AUGUSTUS THUANUS, was a distinguished French historian, and was born at Paris on the 8th of October 1553. In his infancy he was very weak and sickly, and infirmity of health clung to him into man's years. Indisposed to play like other boys by reason of his frail health, he cultivated a taste for drawing and reading when he should have been engaged in the boyish amusements peculiar to his years. De Thou was originally designed for the church, but he allowed his mind to take a much wider course in its thirst for knowledge than was fairly consistent with an aspirant after priest's orders. He greedily studied literature and science, and contracted so great a liking for Cujas, the celebrated jurist, that he set out for Valence in Dauphiné, and attended his lectures on Papinian. Here he formed the acquaintance of Joseph Scaliger, who remained one of his constant friends throughout life. De Thou returned to Paris in 1572 after a year's absence, to witness the marriage of Henry of Navarre; to see the dead body of Coligny hanging from the gibbet; and to be a spectator of the horrible massacre of St Bartholomew. Next year he went to Italy as an attendant of Paul de Foix; he visited Flanders and Holland in 1576, and in 1578 he accepted, with reluctance, the office of an ecclesiastical councillor of the parlement of Paris. In 1582, while on a visit to Bordeaux, he formed the friendship of honest, garrulous old Montaigne, whose character and genius he has warmly praised. His father, Christophe de Thou, first president of the parlement of Paris, having died during the same year, and both his elder and younger brothers being now dead, it was deemed advisable for De Thou to cease to be an ecclesiastic. He accordingly resigned that profession soon after, and accepted the office of Master of Requests, and was appointed to his uncle's place of président du mortier to the parlement of Paris. The De Thous had always remained firm adherents of royalty, and the most illustrious of them would not disgrace his ancestry. He was admitted a counsellor of state to Henry III. at Chartres, and was engaged in various delicate and important missions for Henry III. during the succeeding years. He visited Blois, Paris, Germany, and Italy, all in the royal cause. The death of Henry reached him in Venice in 1589. De Thou immediately set out for France by way of Switzerland, and was received very graciously by Henry IV. at Châteaudun. Henceforward he became a very useful servant to the crown. He was made keeper of the Royal Library on the death of Bishop Amyot in 1591, and in 1593 he commenced his great work Historia Sui Temporis, on which he had meditated for fifteen years.

From this period until his death his time was much occupied in the composition of his history; and except the arrangement of the Edict of Nantes, in which he had an important share, he engaged very little in public business. The death of Henry in 1610 virtually brought his political life to a termination. Soured with public disappointments, and torn by private griefs, caused mainly by the death of his second wife, he died at Paris on the 7th of May 1617, in his sixty-fourth year. A son of De Thou's by his second wife, François August, fell a sacrifice to the inexorable vengeance of Cardinal Richelieu on the 12th of September 1642.

From 1604, when the first eighteen books of De Thou's History appeared, the author was received with acclamation by all the reading public of Europe, as the first historian of the age. During his life there appeared eighty books of it in all, and the remainder, forming in all one hundred and thirty-eight books, were published in 1620. The space over which his record extends, is from the year 1544 to 1607; from the closing years of Francis I. till nearly the completion of the reign of Henry IV. It has all the liveliness and freshness of a narrative, written by a man who was a principal actor in the scenes which it is his business to chronicle; it has, besides, the rare merit of being singularly impartial, the author recognising with uncommon frankness the excellence and even occasional superiority of statesmen and citizens who were avowed rivals to the nation which it was both his duty and his pride to honour and exalt. But it is much more upon the excellence of the workmanship that De Thou's History lives, than upon the comparative accuracy and faithfulness of its recorded facts. While the style of the work certainly displays abundance of easy, flowing eloquence, it is perhaps deficient in the higher quality of vivid picturesqueness. The language in which it was written was Latin; and the author has spared neither time nor labour to make it severely classical. That it does not come quite up to this ideal is, perhaps, not to be wondered at. The best English edition of De Thou's History is that in 7 vols. by Samuel Buckley, 1733. De Thou likewise wrote a number of Latin poems, one of which, entitled De Re Accipitraria (On Hawking), appeared in 1584. The Life of this author has been frequently written. The above notice has been gathered from his autobiography, which terminates in 1601, and from Collinson's Life of Thuanus, London, 1807.

\subsection*{THOUROUT (E0057)}
THOUROUT, a town of Belgium, capital of a canton in the province of West Flanders, 11 miles S.W. of Bruges. It has a fine church, a town-hall, several schools, and benevolent institutions. Linen, hats, leather, cordage, earthenware, soap, tobacco, and other goods are made here. In the vicinity are the remains of an old castle, once the residence of the Counts of Flanders. Pop. 8405.

\subsection*{THRACE (E0058)}
THRACE, the ancient name of a country in the S.E. of Europe, now forming part of European Turkey. In the earliest times the name was used in a very vague sense, and applied to the whole region north of the Aegean Sea, including Macedonia and even Scythia; but when this country became better known to the Greeks, Thrace was confined within narrower and more definite limits. It was bounded on the N. by the Danube, E. by the Euxine and the Bosphorus, S. by the Propontis and the Aegean, and W. by the Strymon, or, at a later period, by the Nestus, which separated it from Macedonia. After the Roman conquest, its size was still further reduced by the chain of Haemus being made its northern boundary. It corresponds partly with the modern Rumelia, which also includes part of Macedonia and Thessaly. (See Rumelia.)

\subsection*{THRAPSTON (E0059)}
THRAPSTON, a market-town of England, in the county and 18 miles N.N.E. of Northampton, on the Nen, which is here crossed by a bridge. The parish church is a cruciform building with a tower and spire; and there are


\section*{Source record}
\begin{description}
\item[Scan ID] \texttt{britannica:309f765b1ecc83e0a7f7762e44b4a927a7555443:train:604}
\item[Image file] \texttt{page-0604-c50a825dd732.jpg}
\item[Image SHA-256] \texttt{c50a825dd732a96ef099cb60b5160d5fadaf7974163160bf5a2ac4446c70b631}
\item[Dataset] NationalLibraryOfScotland/encyclopaedia\_britannica\_illustrated at revision \texttt{309f765b1ecc83e0a7f7762e44b4a927a7555443} (CC0-1.0)
\item[Extraction] model \texttt{nvidia/nemotron-3.5-parse}, raw reply \texttt{extract/raw/page-0604-c50a825dd732.json}
\item[Raw reply SHA-256] \texttt{2bc470de25ef4acb1074a5cd80078ecfb25845f4281841e6bfced9105fa7dad2}
\item[Correction] model \texttt{nvidia/Nemotron-3.5-Super-VL}, block crops, raw replies \texttt{correct/raw/crop\_nothink/page-0604-c50a825dd732-b*.json}; super\_vl: 9
\item[Running header] T H O T H R
\end{description}
\end{document}
